\subsection{Incidences entre moitiés d'émissions régionales}
\label{sec:rec27-empalmes}

Soit $\mathcal U_6$ l'image des six fenêtres consécutives de l'émetteur de \ref{act21:tpk:emision}. Pour $u=(u_1,\ldots,u_6)$, nous définissons $A(u)=(u_1,u_2,u_3)$ et $B(u)=(u_4,u_5,u_6)$. Le graphe des moitiés a pour sommets ces triplets et une arête non orientée entre $A(u)$ et $B(u)$ pour chaque émission.

\begin{proposition}[Cycle des raccordements régionaux]
Les sommets
\[
\begin{aligned}
A&=(830,943,830),&B&=(109,252,252),\\
C&=(810,923,870),&D&=(169,272,272),\\
E&=(850,963,890),&F&=(149,252,212)
\end{aligned}
\]
forment le cycle $A-B-C-D-E-F-A$.
\end{proposition}
\begin{proof}
Nous écrivons chaque germe sous la forme $(i,j,d)$, avec des positions résiduelles entre
$0$ et $8$ et des directions $N=(-1,0)$, $E=(0,1)$ et $S=(1,0)$.
Les couples de germes suivants produisent, au moyen de la récurrence de
54 instants, les six émissions indiquées :
\[
\begin{array}{c|c|c}
\text{émission}&s_+&s_\times\\\hline
A\mid B&(0,6,E)&(0,6,N)\\
B\mid C&(0,6,N)&(1,5,E)\\
C\mid D&(0,6,E)&(0,1,S)\\
D\mid E&(0,6,N)&(0,4,S)\\
E\mid F&(0,6,E)&(0,3,E)\\
F\mid A&(0,6,N)&(0,1,S)
\end{array}
\]
Chaque ligne est un témoin d'appartenance à $\mathcal U_6$, obtenu par
évaluation successive des deux accumulateurs. En séparant chaque émission en
ses deux moitiés, on obtient exactement les six arêtes du cycle.
\end{proof}

Dans la classification de \eqref{eq:palabras-regionales}, \(A\mid B\), \(C\mid D\) et \(E\mid F\) produisent respectivement les mots \(w^{\rm au}\), \(w^{\rm cl}\) et \(w^{\rm pr}\). Les lecteurs régionaux identifient ensuite ces sorties à \(\varphi\), une région de \(\pi\) et \(e\). Les trois autres émissions explicitent leur raccordement dans le graphe. Le résultat décrit des incidences entre sorties de l'émetteur. La réalisation temporelle d'une incidence utilise en outre l'orientation, la phase et la mémoire de l'état enrichi. L'incorporation de l'ancrage APP et des cinq régions de $\pi$ impose des conditions de parcours supplémentaires ; le cycle précédent conserve sa fonction de témoin local de raccordement.

Le cycle de six arêtes appartient à ce graphe non orienté de
demi-émissions. Un parcours temporel enrichi conserve en outre les
données d'ancrage, d'orientation et de fermeture de sa construction. L'incidence
entre deux moitiés et le transport d'un état ont ainsi des domaines et des
conditions de composition explicites.
